%% 
%% The (short) Qpack2 manual
%%

\documentclass[a4paper,11pt]{article}


\usepackage{verbatim}

%- Page layout
%
\pagestyle{empty}
\usepackage[top=30mm,left=30mm,right=30mm,bottom=30mm,nohead]{geometry}


\title{Qpack2}
\author{Patrick Eriksson\\(\texttt{patrick.eriksson@chalmers.se})}
\date{\today}


\newcommand{\newfield}{\vspace{1ex}\noindent}


%--------------------------------------------------------------------------------
\begin{document}


\maketitle

\section{Scope}
\label{sec:scope}
%---
Qpack2 is a package of Matlab functions for performing inversion
work. Only cases where the aim is to retrieve vertical profiles of
atmospheric quantities (1D cases) are handled. The package is also
developed for ground-based observations and inversion of individual
satelitte spectra/scans. This means that inversion of satellite
scanning sequencies where atmospheric fields to be retrieved have a
latitude or longitude extension (2D or 3D) can not be treated by
Qpack2.

The package should be highly general and flexible for covered
measurements. For example, there is basically no limitations regarding
observation geometry. Further, Qpack2 has been developed to be
suitable for operational inversions, as long as not extreme
calculation speed is required. Some important aspects here are that
batch calculations are allowed and atmospheric a priori profiles can
automatically be extracted from climatology data.

The development of Qpack2 is done using the so called optimal
estimation method as the core inversion methodology. Presently no
other inversion methods are handled, but later additions are possible.



\section{Background}
\label{sec:background}
%---
Qpack2 is part of the more general Atmlab package of Matlab functions.
The computational engine of Qpack2 is ARTS-2. The communication with
ARTS-2 is made through the qarts part of Atmlab. OEM inversions are
performed by the function \verb|oem.m| (part of Atmlab).

This can be seen as a direct successor of Qpack, despite all code is
written from scratch. However, for both Qpack versions the
computations are controlled by a set of settings fields packed into a
structure denoted as Q. Qpack was built around ARTS-1, this new Qpack
version is mainly a consequence of that ARTS-2 is the maintained
version. This gave also an oppertunity to make a more stringent
implementation and to add some features. This should hopefully make it
easier to add extensions to Qpack2. Data formats are now more clearly
defined and the ambition level around documentation is higher (but
still relatively modest). A main improvement compared to Qpack2 is
that extraction of a priori data from climatology data is now an
integrated part. On the other hand, the feature of classifying errors
into different categories (0-3 in Qpack) and the associated plotting
features are removed as it was hardly used.

Quality checks are not handled by Qpack2. For example, there is no
defualt check that the frequencies of the measurements are as
expected. This for efficiency reasons and the fact that such checks
are hard to make in a totally general manner. Some checks can be
performed by accompanying functions, such as \verb|qp2_check_f|.


\section{Documentation}
\label{sec:doc}
%---

High level documentation is found in this document. Information on
individual functions is obtained by the standard Matlab help
command. For example:
\begin{verbatim}
>> help qp2_oem
\end{verbatim}
Some of the used data structures are documented following the
pattern expected of the function \verb|qinfo|. For example, to list
all fields (with documentation text) of \verb|Q| (see below):
\begin{verbatim}
>> qinfo(@qpack2)
\end{verbatim}
The documentation for a specific field is obtained as
\begin{verbatim}
>> qinfo(@qpack2,'F_BACKEND')
\end{verbatim}
Wildcards are allowed when naming fields:
\begin{verbatim}
>> qinfo(@qpack2,'ABS_*')
\end{verbatim}


\section{Main data structures}

\subsection{Atmlab settings}
%---
At least the following atmlab settings are used: \verb|ARTS_PATH|,
\verb|ARTS_INCLUDES|, \verb|EXTRAP|, \verb|FMODEL_VERBOSITY|,
\verb|VERBOSITY| and \verb|WORK_AREA|.

For further information:
\begin{verbatim}
>> arts('-d z_surface');
\end{verbatim}


\subsection{Forward model and retrieval settings}
\label{sec:Q}
%---
Forward model parameters and retrieval variables are joined into a
single structure, denoted as \verb|Q|. This structure has a fixed set
of fields. That is, the structure must always contain exactly this set
of fields. However, some fields are used only in particularly
circumstances and not all fields must be set.

The fields of \verb|Q| are listed and described in
Appendix~\ref{app:q}. The same information is obtained inside Matlab
through \verb|qinfo|, where the defintion function is
\verb|qpack2|. The examples in Section~\ref{sec:doc} are here valid.

Many of the fields of \verb|Q| are identical to the corresponding
field of the qarts interface, and additional information can be found
as:
\begin{verbatim}
>> qinfo(@qarts)
\end{verbatim}
On-line documentation of ARTS methods and workspace variables is
obtained in a terminal as
\begin{verbatim}
$ arts -d z_surface
\end{verbatim}
or from within Matlab as
\begin{verbatim}
>> arts('-d z_surface');
\end{verbatim}


\subsection{Measurement data}
\label{sec:Y}
%---
The series of measurements (spectra or scans) to be treated is packed
into a data structure denoted as \verb|Y|. The data fields of this
structure are described in \verb|qp2_y|, and available information can be
obtained as
\begin{verbatim}
>> qinfo(@qp2_y)
\end{verbatim}
The same information is appended to this document in
Appendix~\ref{app:y}. An example:
\begin{verbatim}
>> Y(3)

ans = 
          DAY: 3
            F: {}
         HOUR: 11
        HSE_P: 10000
        HSE_Z: 16000
     LATITUDE: {}
    LONGITUDE: {}
       MINUTE: 57
        MONTH: 3
       SECOND: 58.0000
       TNOISE: 0.0243
            Y: [800x1 double]
         YEAR: 2002
           ZA: 65
\end{verbatim}



\subsection{Retrieved data}
\label{sec:P}
%---
Retrieved data are returned as a structure array, in the same manner
as the measurements are provided (\verb|Y|). The returned fields
depend on the selection of retrieval quantities and on user
settings. This makes it impossible to use qinfo here (as for \verb|Q|
and \verb|Y|), and the fields are instead described below in this
document.

An example on a basically minimal \verb|P|:
\begin{verbatim}
K>> P(1)

ans = 
             cost: 1.3355
           cost_x: 0.0021
           cost_y: 1.3334
             year: 2002
            month: 3
              day: 1
             hour: 20
           minute: 28
           second: 28.0000
         latitude: 58
        longitude: -12
           tnoise: 0.0274
        converged: NaN
    species1_name: 'H2O'
       species1_p: [50x1 double]
       species1_x: [50x1 double]
      species1_xa: [50x1 double]
       species1_e: [50x1 double]
      species1_eo: [50x1 double]
      species1_es: [50x1 double]
      species1_mr: [50x1 double]
\end{verbatim}



\subsubsection{Fields always returned}
%
The following fields are always provided:

\newfield 
\verb|converged|:\\
Convergence flag. See further \verb|oem.m|.

\newfield 
\verb|cost|, \verb|cost_x| and \verb|cost_y|:\\
Cost values of retrieved state and fit to measurement. See further \verb|oem.m|.

\newfield 
\verb|P.latitude| and \verb|P.longitude|:\\
The geographical position of the measurement.

\newfield 
\verb|tnoise|:\\
Mean of assumed thermal noise. This is simple average, no weighting
with channel widths is applied. This value has no direct physical
meaning, but can be used as a mean to e.g. remove the most noisy
measurements (without having \verb|Y| at hand).


\subsubsection{With Q.OUT\_Y  true}
%
The following fields are also provided if \verb|Q.OUT_Y| is set to true:

\newfield 
\verb|f|:\\
Frequency vector of measurement. Equals \verb|Y.F| if set, otherwise
\verb|Q.F_BACKEND|.

\newfield 
\verb|y|:\\
Measured spectrum.

\newfield 
\verb|yf|:\\
Fitted spectrum. That is, the spectrum matching retrieved state vector.

\newfield
\verb|baseline|:\\
Retrieved ``baseline''. That is, retrieved spectrum distortion through
e.g.\ polynomial fitting.


\subsubsection{With Q.OUT\_PTZ  true}
%
The following fields are also provided if \verb|Q.OUT_PTZ| is set to true:

\newfield 
\verb|p|, \verb|t| and \verb|z|:\\
A consistent set of pressure, temperature and altitude vectors. The
vectors are identical to the ones that are provided to the forward
model. If temperatures are retrieved, this information is used.
Otherwise returned data match climatology temperatures. Altitudes are
obtained by applying hydrostatic equilibrium.


\subsubsection{Species fields}
%
The retrieved species are indexed after the order they are specified
in \verb|Q| and named as \verb|species1|, \verb|species2| \dots A
number of fields is returned for each species. The field names below
are valid for species 1. If two species are retrieved there is also a
field named as \verb|species2_name| etc. The following fields
are always returned:

\newfield 
\verb|species1_name|:\\
Tag name of species.

\newfield 
\verb|species1_p|:\\
Vertical/pressure retrieval grid for the species.

\newfield 
\verb|species1_x|:\\
Retrieved profile. Unit depends on \verb|Q.SPECIES(i).UNIT|.

\newfield 
\verb|species1_xa|:\\
A priori state for the species.

\newfield 
\verb|species1_e|:\\
Total retrieval (observation + smooting) error for the species.

\newfield 
\verb|species1_eo|:\\
Observation error for the species.

\newfield 
\verb|species1_es|:\\
Smoothing error for the species.

\newfield
\verb|species1_mr|:\\
Measurement response for the species. This is the sum of the rows of
\verb|species_A| (see below).


\subsubsection{Additional species fields}
%
If \verb|Q.OUT_A| is set to be true, the following field is also
included:

\newfield
\verb|species1_A|:\\
The species specific averaging kernels. That is, this is the diagonal
centred sub-matrix covering only the resolution between changes of the
species at different altitudes. The resolution between variables of
different retrieval quantities is thus not covered.


\subsubsection{Baseline fit variables}
%
If any baseline fit is performed and \verb|Q.OUT_BASELINE| is set to true,
the following field is also provided

\newfield
\verb|P.polyfit0_x|, \verb|P.polyfit0_xa|, \verb|P.polyfit0_e|,
\verb|P.polyfit0_eo|, \verb|P.polyfit0_es| and \verb|P.polyfit0_mr|:\\







\section{Standard inversions}
\label{sec:stdinversions}
%---


\verb|Q.F_BACKEND|

\begin{verbatim}
>> Q = q_h2o;
>> Y = get_onsala_Y( 'h2o', date1, date2 );
>> P = qp2_oem( Q, Y );
\end{verbatim}


\begin{verbatim}
>> if ~qp2_check_f( Q, Y, 1e3 );
>>   error( 'Some mismatch between Q.F_BACKEND and frequencies of spectra.' );
>> end
>> [Y.F] = deal( {} ); % The field F is now obselete. {} is used inside 
>>                     % qarts and qpack2 to flag undefined fields
\end{verbatim}



\appendix

\newpage\noindent{\huge \bf Appendices}

\section{The fields of Q}
\label{app:q}
%
\verbatiminput{qpack2.txt}


\newpage
\section{The fields of Y}
\label{app:y}
%
\verbatiminput{qp2_y.txt}



\end{document}


#############################################################################
###   Qpack2
#############################################################################

Scope
-
Qpack2 shall handle:
1. Batch inversions of 1D cases (only single spectra cases, but handling
"double" observation modes such as beam switching and frequency switching).
2. Produce spectra matching the a priori assumptions.

Option 2 is triggered by setting all spectra (Y.y) to be empty.

Quality checks are not hanlded by qpack2. For example, there is no
check that the frequencies of the measurements are as expected
(Q.F_BACKEND). This for efficiency reasons and the fact that such
checks are hard to make in a totally general manner. Some checks can
be performed by accompanying functions, such as qp2_check_f. 

Beside the settings below, some general atmlab settings are used. This 
includes
ARTS_PATH
ARTS_INCLUDES
EXTRAP
FMODEL_VERBOSITY
VERBOSITY
WORK_AREA

A calculation run could like this:

>> Q = q_h2o;
>> Y = get_onsala_Y( 'h2o', date1, date2 );
>> qp2_check_f( Q, Y, 1e3 );
>> Y = rmfield( Y, 'f' );
>> P = qp2_oem( Q, Y );





#############################################################################
### Y : Measurement data
#############################################################################

The measurements are provided to qpack2 as an array of structures. 
Mandatory fields are:

Y.y
-
The spectrum. All spectra in the array must have the same length.

Y.year
Y.month
Y.day
-
The time of the measurement. All three fields are numeric. This
information is primarily used to extract data from the climatology
database provided.

Y.zenith_angle
- 
Line of-sight zenith angle for the measurement. A scalar value.

Y.tnoise_si
-
Magnitude of thermal noise (1 std). A scalar or a vector having the
same length as y. If a scalar, the value is applied for all
spectrometer channels.

Y.hse_p
Y.hse_z
-
The reference point when enforcing hydrostatic equilibrium. The
geometrical altitude (hse_z) is given for one presssure (hse_p).


Optional fields are:

Y.f
-
Frequency for each y value. Selected frequency should match Q.F_BACKEND.

Y.hour
Y.minute
Y.second
-
Allowing a more detailed specification of measurement time.

Y.latitude
Y.longitude
-
The geographical position of the measurement. If not given, the 
position is taken from Q.PLATFORM_LAT and Q.PLATFORM_LON.

An example:

>> Y(3)

ans = 

            year: 2002
           month: 3
             day: 3
            hour: 11
          minute: 57
          second: 58.0000
    zenith_angle: 65
           hse_p: 10000
           hse_z: 16000
       tnoise_si: 0.0243
               y: [800x1 double]





#############################################################################
### Q : Forward model and retrieval data
#############################################################################

All fields are described and set to default/dummy values in
qpack2.m. Use qinfo to obtain the information. For everything:

qinfo( @qpack2 )

To read about a specific field, such as POLYFIT_SI:

qinfo( @qpack2, 'POLYFIT_SI' )





#############################################################################
### P : Result and auxiliary data of retrieval
#############################################################################

A structure array. The following fields are always provided
---

P.latitude
P.longitude
-
Geographical position.

P.year
P.month
P.day
P.hour
P.minute
P.second
-
Measurement time. 

P.tnoise_si
-
Mean of assumed thermal noise (a direct averaging, no weighting with
channel widths). This value has no direct physical meaning, but can be
used as a mean to e.g. remove the most noisy measurements (without 
having Y at hand).

P.converged
-
See oem.m.

P.cost
P.cost_y
P.cost_x
-
See oem.m.


The following fields are also provided if Q.OUT_Y is set to true:
---

P.f
-
Frequency vector of measurement (Q.F_BACKEND).

P.y
-
Measured spectrum.

P.yf
-
Fitted spectrum. That is, the spectrum matching retrieved state vector.


The following fields are also provided if Q.OUT_PTZ is set to true:
---

P.p P.t P.z 
-
A consistent set of pressure, temperature and altitude vectors. The
vectors are identical to the ones that are provided to the forward
model. If temperatures are retrieved, this information is used.
Otherwise returned data match climatology temperatures. Altitudes are
obtained by applying hydrostatic equilibrium.


If any species is retrieved, these fields are always returned:
---

species1_name
-
Tag name of species.

species1_p
-
Vertical/pressure retrieval grid for the species.

species1_x
- 
Retrieved profile. Unit depends on Q.SPECIES(i).UNIT.

species1_xa
-
A priori state for the species.

species1_e
-
Total retrieval (observation + smooting) error for the species.

species1_eo
-
Observation error for the species.

species1_eo
-
Smoothing error for the species.

species1_mr
-
Measurement response for the species.



If Q.OUT_A is set to true, the following field is also provided
---

P.species1_A
-
Averaging kernels for the species. 



If more than one species is retrieved, data for other species are
returned as species2_name etc.



If any baseline fit is performed and Q.OUT_BASELINE is set to true,
the following field is also provided

P.baseline
-
Retrieved baseline fit.



Result for each baseline fit variable is also provided. For example,
the following fields are added to handle polynomial coefficent 0:

P.polyfit0_x
P.polyfit0_xa
P.polyfit0_e
P.polyfit0_eo
P.polyfit0_es
P.polyfit0_mr
